
\begin{figure}[H]
\centering
% \begin{adjustbox}{center,max width=\textwidth}

\resizebox{\textwidth}{!}{
\begin{tikzpicture}[
    box/.style={rectangle,draw,thick,rounded corners,minimum height=0.8cm},
    process/.style={box,fill=gray!20},
    neuro/.style={box,fill=green!20},
    behavior/.style={box,fill=orange!20},
    env/.style={box,fill=teal!15,minimum width=3.6cm,minimum height=1.2cm},
    arrow/.style={->,>=Stealth,thick},
    m2env/.style={arrow,draw=orange!80},        % Mother -> Environment
    env2m/.style={arrow,draw=teal!70!black},    % Environment -> Mother
    c2env/.style={arrow,draw=blue!70},          % Child -> Environment
    env2c/.style={arrow,draw=violet!80!black,line width=1pt,
    preaction={draw=white,line width=3pt}}, % Environment -> Child
    feedback/.style={arrow,draw=blue!50,dashed} % internal feedback
]

% ---------- KEY: symmetric bounding box ----------
\path[use as bounding box] (-9,0) rectangle (9,9.6);

%% MOTHER SYSTEM (Left side - Full model)
\node[font=\bfseries] at (-5,9.5) {MOTHER AGENT - Full Model};

% Percept -> Stimuli
\node[process] (m_percept) at (-7.3,7.5) {Perception};
\node[process] (m_stimuli) at (-4.39,7.5) {Stimuli};
\draw[arrow] (m_percept) -- (m_stimuli);

% Neuroendocrine
\node[neuro,minimum width=3cm] (m_neuro) at (-3.7,5.5) {Neuroendocrine};
% \node[font=\tiny] at (-5,5.6) {OT, AVP, DA};
\draw[arrow] (m_stimuli.south east) -- (m_neuro.north);

% Biological process
\node[neuro,minimum width=3cm] (m_bio) at (-7.5,5.5) {Biological Processes};
% \node[font=\tiny] at (-5,5.6) {OT, AVP, DA};
\draw[arrow] (m_stimuli.south west) -- (m_bio.north);

% Motivation
\node[process] (m_motiv) at (-5.3,3.5) {Motivation};
\draw[arrow] (m_bio.south) to (m_motiv);
\draw[arrow] (m_neuro.south) to  (m_motiv);


% Selected Action
\node[behavior] (m_behave) at (-5.3, 1.5) {Selected Behavior};
% \node[font=\tiny] at (-5,0.9) {Approach, Feed, Protect, Warm};
\draw[arrow] (m_motiv) to [out=-90, in=90] (m_behave);

\draw[feedback] (m_behave.west) to[out=180,in=200]
  node[pos=0.5,left,font=\tiny,xshift=55pt,yshift=2pt]
  {Behavior effect} (m_bio.south west);
\draw[feedback] (m_bio.north west) to[out=150,in=210]
   node[pos=0.5,left,font=\tiny,xshift=65pt,yshift=-5pt]
   {Allostatic regulation} (m_percept.west);

% %% CHILD SYSTEM (Right side - Simplified model)
\node[font=\bfseries] at (4.5,9.5) {CHILD AGENT - Simplified Model};

%Perception
\node[process,fill=blue!20] (c_percept) at (3,7.5) {Limited Perception};

% Stimuli
\node[process,fill=blue!20] (c_stimuli) at (6.7,7.5) {Stimuli};
\draw[arrow] (c_percept) -- (c_stimuli);

% Internal States
\node[neuro,fill=blue!20] (c_needs) at (5,5.5) {Internal States};
\draw[arrow] (c_stimuli) -- (c_needs);

% Need Level (derived from internal states)
\node[process,fill=blue!20] (c_distress) at (5,3.5) {Need Level};
\draw[arrow] (c_needs) -- (c_distress) ;

% In implementation, the mother directly observes child state variables;
% we depict this as a need signal computed from internal states.
\node[behavior,fill=blue!20] (c_signal) at (5,1.5) {Need Signals};
\draw[arrow] (c_distress) -- (c_signal);

\draw[feedback] (c_signal.east) to[out=0,in=0]
  node[pos=0.2,left,font=\tiny,xshift=60pt,yshift=2pt]
 {Comfort received} (c_needs.east);






%% ENVIRONMENT (center)
\node[env] (env) at (0,5.5) {Environment};

% Mother <-> Environment
\draw[m2env] (m_behave.east) to[out=0,in=200]  
 node[pos=0.5,left,font=\tiny,xshift=50pt,yshift=2pt]
 {Action to env.} (env.south west);
\draw[env2m] (env.north) to[out=90,in=30] 
 node[pos=0.5,left,font=\tiny,xshift=8pt,yshift=-5pt]
 {Resources and Threats cue}(m_percept.north);

%% Background highlight boxes (requires backgrounds library)
\begin{scope}[on background layer]
    \node[rectangle,draw=red!50,thick,dashed,minimum width=10cm,minimum height=9.5cm, fill=red!5] at (-5,5.4) {};
    \node[rectangle,draw=blue!50,thick,dashed,minimum width=9cm,minimum height=9.5cm, fill=blue!5] at (4.5,5.4) {};
\end{scope}

% %% Child <-> Environment
\draw[c2env] (c_signal.west) to[out=180, in=340]
node[pos=0.5,left,font=\tiny,xshift=0pt,yshift=0pt]
{Need cues}(env.south east);

\draw[env2c] (env.north) to[out=90,in=180]
node[pos=0.4,left,font=\tiny,xshift=35pt,yshift=-5pt]
{Safety cue}(c_percept.west);

% Title
\node[font=\Large\bfseries] at (0, 10.8)
  {Overview Neuroendocrine Systems};

\end{tikzpicture}
}
\caption{Overview of the proposed neuroendocrine-inspired agent architecture. 
The mother agent integrates perception, biological processes, and neuroendocrine modulation to generate motivations and select behaviors. 
The child agent is simplified and represented by internal need states; these provide need cues that influence maternal perception and motivation.}
\label{fig:neuroendocrine_systems}
\end{figure}

